% Source: https://artofproblemsolving.com/community/c3h2270974_sequences
% Stable identity: aops-topic-2270974-q01
The radii are $8,4,2,1,\ldots$ cm, so the areas are $64\pi,16\pi,4\pi,\pi,\ldots$: each area is $\dfrac{1}{4}$ of the previous one (halving a length divides an area by $4$).

\begin{asy}
usepackage("amsmath");
size(7cm);
defaultpen(fontsize(10pt));
real x=0; real d=16;
for(int i=0;i<6;++i){ draw(circle((x+d/2,0),d/2)); if(i<4) label("$"+string(d)+"$",(x+d/2,0)); x+=d; d/=2; }
label("$\cdots$",(x+0.8,0));
\end{asy}

This is an infinite geometric series with first term $64\pi$ and ratio $\dfrac{1}{4}$, so its sum is
\[\dfrac{64\pi}{1-\dfrac{1}{4}}=\dfrac{64\pi}{\dfrac{3}{4}}=\boxed{\dfrac{256\pi}{3}}\text{ cm}^2\approx268\text{ cm}^2.\]
